\section{Discussion}
\label{sec:discussion}

The finding that TIGER profiles predict BGG ratings at $r \geq 0.60$ (versus $r \approx 0.40$
for weight alone) establishes that design character is not independent of player satisfaction.
Games that score high on Elegance and Game-Architecture are better liked, at least in
the BGG community, independently of their complexity level. This is a non-trivial finding:
it means that optimizing for TIGER profile dimensions is not purely an aesthetic exercise
but has measurable consequences for the player experience outcomes that drive long-term
game ratings.

\subsection{Implications for Pre-Release Quality Forecasting}

A publisher evaluating a prototype can estimate TIGER scores from the design document and
a small number of early playtest sessions, before a full playtest dataset is available.
The regression model can then provide a directional rating forecast: not a precise point
estimate, but a signal about whether the prototype's design character is in a region of
TIGER space associated with high long-term ratings. The main epistemic challenge is that
early TIGER estimates carry scoring uncertainty that propagates into the rating forecast.
Based on the inter-rater reliability data from TG-A1~\cite{dellaLibera2026corpus},
Tension and Elegance scores stabilize at $\alpha \geq 0.85$ after 3--5 playtests, while
Game-Architecture and Range require 8--10 sessions. A practical recommendation is to
compute the forecast twice: once with early T and E estimates (the two highest-impact
ablation contributors) and once with full five-dimension scores, using the delta as a
measure of residual forecast uncertainty. This two-stage approach preserves directional
validity while communicating honest confidence bounds to decision-makers.

\subsection{Limitations}

BGG ratings are community-derived and reflect the preferences of a self-selected population
(predominantly Western, predominantly male, predominantly hobbyist). The regression model
predicts satisfaction for this population, not for general consumer or family audiences.
Extending the prediction target to other satisfaction measures (retail sales, Spiel des
Jahres nominations) is future work. The held-out test set ($N = 28$) is small enough
that point estimates of $r$ and RMSE carry substantial uncertainty; confidence intervals
will be reported in the final version.
